\documentclass[11pt]{article}
\usepackage[margin=1in]{geometry}
\usepackage{amsmath,amssymb,amsthm}
\usepackage{hyperref}
\usepackage{enumitem}

\newtheorem{theorem}{Theorem}
\newtheorem{lemma}[theorem]{Lemma}
\newtheorem{corollary}[theorem]{Corollary}
\theoremstyle{definition}
\newtheorem{definition}[theorem]{Definition}
\theoremstyle{remark}
\newtheorem{remark}[theorem]{Remark}

\usepackage{xurl}
\let\oldus\_
\renewcommand{\_}{\oldus\allowbreak}
\newcommand{\Per}{\operatorname{Per}}
\newcommand{\Fix}{\operatorname{Fix}}

\title{A disproof of Conjecture 00000000741\\[2pt]
\large The squaring map never has exactly three periodic points}
\date{}

\begin{document}
\sloppy
\maketitle

\section{The conjecture}

Conjecture 00000000741 reads:

\begin{quote}
\emph{Definition:} Let $n(p)$ denote the number of periodic points (counting distinct periodic orbits) of
$f(x) = x^2$ on $\mathbb{Z}_p$. \emph{Conjecture:} For a density-one set of primes, $n(p) = 3$ (the fixed
points $0$, $1$ and $-1$); the exceptional primes (with extra periodic points) have count
$\le x^{1/2+\varepsilon}$, and all exceptions are characterized by $p^2 - 1$ containing large power factors.
\end{quote}

The Chinese version defines $n(p)$ as the number of periodic points and specifies in its parenthesis that
the count is of the distinct points lying on periodic orbits. We therefore read $n(p)$ as the number of
periodic points of $f$. For $\mathbb{Z}_p$ we treat both usual meanings: the field $\mathbb{Z}/p\mathbb{Z}$
and the ring of $p$-adic integers. The argument works in every integral domain, so both are covered at once.

The conjecture is a conjunction of three clauses:
\begin{enumerate}[label=(\roman*)]
  \item (\emph{density one}) there is a set $S$ of primes of relative density one among the primes such that
        $n(p) = 3$ for every $p \in S$;
  \item (\emph{exception count}) the number of exceptional primes up to $x$ is at most $x^{1/2+\varepsilon}$;
  \item (\emph{characterization}) all exceptions are characterized by $p^2-1$ having large power factors.
\end{enumerate}
We show that clause (i) is false. Hence the conjunction (i)$\wedge$(ii)$\wedge$(iii) is false, whatever
(ii) and (iii) mean. In fact $n(p) \ne 3$ for \emph{every} prime $p$, so the set of primes with $n(p)=3$ is
empty.

\section{Definitions}

\begin{definition}[Periodic point]
Let $f : X \to X$. A point $x$ is \emph{periodic} if $f^{k}(x) = x$ for some $k \ge 1$, where $f^k$ is the
$k$-fold iterate. It is \emph{fixed} if $f(x) = x$. We write $\Per(f)$ and $\Fix(f)$ for these sets. This is
Mathlib's \texttt{Function.periodicPts f}, the set of $x$ such that \texttt{IsPeriodicPt f n x}
($f^{n}(x) = x$) for some $n > 0$, and \texttt{Function.fixedPoints f}.
\end{definition}

\begin{definition}[The count $n(p)$]
For a ring $R$ let $f_R(x) = x^2$ and $N(R) = |\Per(f_R)|$ (Lean: \texttt{numPeriodic R}, the
\texttt{Nat.card} of the set). Then $n(p) = N(\mathbb{Z}/p\mathbb{Z})$ (Lean: \texttt{n p}) and, for $p$
prime, $n_{\mathrm{adic}}(p) = N(\mathbb{Z}_p)$ (Lean: \texttt{nPadic p}).
\end{definition}

\begin{definition}[Relative density among the primes]
For $S \subseteq \mathbb{N}$ let
\[
  r_S(N) = \frac{\#\{p < N : p \text{ prime}, \ p \in S\}}{\#\{p < N : p \text{ prime}\}}
\]
(Lean: \texttt{primeRatio S N}; the denominator is Mathlib's \texttt{Nat.primeCounting' N}). $S$ has
relative density one if $r_S(N) \to 1$ as $N \to \infty$ (Lean: \texttt{HasPrimeDensityOne S}). Clause (i)
for a count function $m$ is \texttt{DensityOneClause m}: some $S$ with $r_S(N)\to 1$ has $m(p) = 3$ for
every prime $p\in S$.
\end{definition}

\section{The proof}

Throughout, $R$ is an integral domain and $f(x) = x^2$. Since $f^{k}(x) = x^{2^k}$ (Lean:
\texttt{iterate\_sq}), $x \in \Per(f)$ if and only if $x^{2^k} = x$ for some $k \ge 1$
(\texttt{mem\_periodicPts\_sq}). Clearly $0, 1 \in \Per(f)$.

\begin{lemma}[\texttt{fixedPoints\_sq}]\label{lem:fix}
$\Fix(f) = \{0, 1\}$.
\end{lemma}
\begin{proof}
$x^2 = x$ means $x(x-1) = 0$, so $x = 0$ or $x = 1$ because $R$ has no zero divisors.
\end{proof}

\begin{lemma}[\texttt{neg\_one\_mem\_periodicPts\_iff}]\label{lem:neg}
$-1 \in \Per(f)$ if and only if $-1 = 1$ in $R$.
\end{lemma}
\begin{proof}
For $k \ge 1$ the exponent $2^k$ is even, so $(-1)^{2^k} = 1$; thus $(-1)^{2^k} = -1$ forces $1 = -1$.
Conversely, if $-1 = 1$ then $-1$ is fixed.
\end{proof}

\begin{lemma}[\texttt{exists\_inv\_periodic}]\label{lem:inv}
If $x \in \Per(f)$ and $x \ne 0$, there is $y \in \Per(f)$ with $xy = 1$.
\end{lemma}
\begin{proof}
Let $x^{2^k} = x$ with $k \ge 1$. Then $x \cdot x^{2^k-1} = x \cdot 1$, and cancelling $x$ gives
$x^{2^k - 1} = 1$. Put $y = x^{2^k-2}$; then $xy = x^{2^k-1} = 1$. Moreover
$x\, y^{2^k} = x^{2^k} y^{2^k} = (xy)^{2^k} = 1 = xy$, and cancelling $x$ gives $y^{2^k} = y$.
\end{proof}

\begin{theorem}[\texttt{card\_periodicPts\_ne\_three}]\label{thm:main}
In every integral domain $R$, the squaring map does not have exactly three periodic points:
$N(R) \ne 3$.
\end{theorem}
\begin{proof}
Suppose $|\Per(f)| = 3$. Then $\Per(f)$ is finite and contains the two distinct points $0, 1$, so it contains
a third point $x \notin \{0,1\}$, and $\Per(f) = \{0, 1, x\}$. By Lemma~\ref{lem:inv} there is
$y \in \Per(f)$ with $xy = 1$. Here $y \ne 0$ because $xy = 1$, and $y \ne 1$ because otherwise $x = 1$.
Hence $y = x$ and $x^2 = 1$, so $(x-1)(x+1) = 0$ and $x = -1$. By Lemma~\ref{lem:neg}, $-1 = 1$, so $x = 1$,
a contradiction.
\end{proof}

\begin{corollary}[\texttt{n\_ne\_three}, \texttt{nPadic\_ne\_three}]
For every prime $p$, $n(p) \ne 3$, both on $\mathbb{Z}/p\mathbb{Z}$ and on the $p$-adic integers.
\end{corollary}
\begin{proof}
Both rings are integral domains (Mathlib provides the instances).
\end{proof}

\begin{theorem}[\texttt{not\_densityOneClause}]
If $m(p) \ne 3$ for every prime $p$, then \texttt{DensityOneClause m} is false. In particular clause (i)
fails for $n$ and for $n_{\mathrm{adic}}$.
\end{theorem}
\begin{proof}
Let $S$ have $m(p) = 3$ for every prime $p \in S$. Then no prime lies in $S$, so the numerator of
$r_S(N)$ is $0$ for every $N$, and $r_S(N) = 0$ for every $N$. The constant sequence $0$ does not converge
to $1$.
\end{proof}

We also record (Lean: \texttt{neg\_one\_not\_periodic}) that for odd primes $p$ the element $-1$ is not
periodic, because $-1 \ne 1$ in $\mathbb{Z}/p\mathbb{Z}$. Together with Lemma~\ref{lem:fix}, this shows that
the three points named in the conjecture, ``the fixed points $0$, $1$ and $-1$'', are never all fixed or
even periodic for $p > 2$.

\begin{remark}[Exact value, not used in the proof]
In $\mathbb{Z}/p\mathbb{Z}$ a unit $x$ is periodic exactly when its multiplicative order is odd, so
$n(p) = 1 + m$, where $m$ is the odd part of $p-1$. Thus $n(p)$ is always even:
$n(2)=n(3)=n(5)=2$, $n(7)=4$, $n(11)=6$, $n(13)=4$. A Python check confirms this formula for all
$p < 1500$. So every prime except $2$ and the Fermat primes has more than three periodic points.
\end{remark}

\begin{remark}[The orbit-count reading]
The English parenthesis ``counting distinct periodic orbits'' could also be read as counting orbits rather
than points. Under that reading the claim is false as well, although this remark is not formalized. For
every prime $p \equiv 1 \pmod 7$, the group $(\mathbb{Z}/p\mathbb{Z})^\times$ contains six elements of
order $7$. Since $2$ has order $3$ modulo $7$, squaring splits them into two $3$-cycles. Together with
$\{0\}$ and $\{1\}$ this gives at least four periodic orbits. By Dirichlet's theorem in its density form,
these primes have relative density $1/6$, so the primes with exactly three orbits do not have density one,
and the exceptions are not $O(x^{1/2+\varepsilon})$. Even under this reading, the three orbits are never
``the fixed points $0$, $1$ and $-1$'', since $-1$ is not fixed for $p>2$ (Lemma~\ref{lem:fix}).
\end{remark}

\section{Lean formalization}

File \texttt{lean/Conjecture741/Basic.lean} (Lean 4.33.1, Mathlib v4.33.1; only \texttt{import Mathlib}).
The main theorem is
\begin{quote}
\texttt{theorem conjecture\_00000000741\_false (Q : Prop) :}\
\quad $\neg(\texttt{DensityOneClause n} \wedge Q) \;\wedge\; \neg(\texttt{DensityOneClause nPadic} \wedge Q)
\;\wedge$\
\quad $(\forall p,\ p\ \text{prime} \to \texttt{n}\ p \ne 3) \;\wedge$\
\quad $(\forall p,\ p\ \text{prime} \to \texttt{fixedPoints}\,(y \mapsto y^2 : \mathbb{Z}/p \to \mathbb{Z}/p)
= \{0,1\}) \;\wedge$\
\quad $(\forall p,\ p\ \text{prime} \to p \ne 2 \to (-1 : \mathbb{Z}/p) \notin
\texttt{periodicPts}\,(y \mapsto y^2))$
\end{quote}
The proposition \texttt{Q} stands for any formalization of clauses (ii) and (iii). The supporting
results are \texttt{iterate\_sq}, \texttt{mem\_periodicPts\_sq}, \texttt{fixedPoints\_sq},
\texttt{neg\_one\_mem\_periodicPts\_iff}, \texttt{exists\_inv\_periodic},
\texttt{card\_periodicPts\_ne\_three}, \texttt{n\_ne\_three}, \texttt{nPadic\_ne\_three},
\texttt{neg\_one\_not\_periodic} and \texttt{not\_densityOneClause}.

\paragraph{Axiom audit.} \verb|#print axioms Conjecture741.conjecture_00000000741_false| reports
\texttt{[propext, Classical.choice, Quot.sound]}. There is no \texttt{sorry} and no
\texttt{native\_decide}, and the file also compiles with \texttt{-DwarningAsError=true}.

\end{document}
